\documentclass[11pt,a4paper]{article}
\usepackage[T1]{fontenc}
\usepackage[utf8]{inputenc}
\usepackage{lmodern}
\usepackage[margin=22mm,headheight=16pt]{geometry}
\usepackage{amsmath,amssymb,graphicx,array,longtable,booktabs}
\usepackage{microtype,xcolor,hyperref,xurl,fancyhdr}
\definecolor{Navy}{HTML}{17324D}
\definecolor{Teal}{HTML}{087F8C}
\definecolor{Gold}{HTML}{B88128}
\hypersetup{colorlinks=true,linkcolor=Navy,urlcolor=Teal,bookmarksnumbered=true,hypertexnames=false}
\setlength{\parindent}{0pt}
\setlength{\parskip}{6pt}
\setlength{\emergencystretch}{4em}
\setcounter{tocdepth}{1}
\pagestyle{fancy}\fancyhf{}
\fancyfoot[R]{\small\thepage}
\fancyfoot[L]{\small\sffamily Sharif University of Technology}
\renewcommand{\headrulewidth}{0.3pt}
\makeatletter
\renewcommand\section{\@startsection{section}{1}{\z@}{-3.2ex plus -1ex minus -.2ex}{1.6ex plus .2ex}{\normalfont\Large\sffamily\bfseries\color{Navy}}}
\renewcommand\subsection{\@startsection{subsection}{2}{\z@}{-2.5ex plus -1ex minus -.2ex}{1.2ex plus .2ex}{\normalfont\large\sffamily\bfseries\color{Teal}}}
\makeatother
\fancyhead[L]{\small\sffamily\color{Navy}Sharif University of Technology}\fancyhead[R]{\small\sffamily\color{Teal}Project Assignment}\begin{document}
\begin{titlepage}\centering\vspace*{7mm}
\includegraphics[width=31mm]{assets/sharif_emblem_en.png}\par\vspace{5mm}
{\Large\sffamily\bfseries\color{Navy}SHARIF UNIVERSITY OF TECHNOLOGY\par}
\vspace{2mm}{\sffamily Department of Electrical Engineering\par}
\vspace{17mm}{\color{Teal}\rule{\linewidth}{1.4pt}}\par\vspace{6mm}
{\large\sffamily\color{Teal}PROJECT ASSIGNMENT\par}
\vspace{8mm}{\huge\sffamily\bfseries\color{Navy}Single-Node Hadoop and Spark Deployment with Distributed Word Count\par}
\vspace{9mm}{\color{Gold}\rule{45mm}{1pt}}\par\vfill
\begin{minipage}{.88\linewidth}
{\sffamily\color{Teal}INSTRUCTOR / INSTRUCTORS}\par
Dr Iman Gholampour\par\vspace{7mm}
{\sffamily\color{Teal}ACADEMIC CONTEXT}\par
Big Data Processing\par\end{minipage}\vfill
\end{titlepage}
\tableofcontents\clearpage
\section{Purpose and submission context}

The first Big Data Analytics assignment, version 1.0, belongs to the first semester of academic year 1404-1405. The instructor is Dr Iman Gholampour;\allowbreak{} the assignment designers are Hossein Soleimani and Ali Habibi. The exercise combines concise conceptual answers with a documented Linux deployment of HDFS and Spark. A single machine may serve as NameNode, DataNode, Spark master, and Spark worker. This is a laboratory deployment, not evidence of a production multi-machine cluster.

\section{Sixteen conceptual questions}

\begin{enumerate}

\item Explain batch and streaming processing in Spark and give appropriate scenarios for each.

\item Define an RDD and identify workloads where it remains preferable to DataFrame or Dataset APIs.

\item Compare DataFrames and Datasets, including advantages, disadvantages, and language constraints.

\item Describe the responsibilities of the driver and executors.

\item Explain shuffle, why it becomes a performance bottleneck, and give an example.

\item Explain when cache/\allowbreak{}persist helps and compare persistence storage levels.

\item Identify causes of executor OutOfMemory and practical remedies.

\item Describe Adaptive Query Execution and the optimizations it provides.

\item Explain Tungsten and whole-stage code generation.

\item Compare Standalone, YARN, and Kubernetes deployment and their appropriate contexts.

\item Explain why Parquet and ORC can outperform CSV, discussing column statistics and compression.

\item Describe Delta Lake capabilities such as ACID transactions, time travel, and Z-ordering.

\item Discuss when a traditional data warehouse or a lakehouse is appropriate.

\item Compare FIFO and Fair scheduling with examples.

\item Explain scheduler pools and how to separate ETL from ad hoc queries.

\item Explain how dynamic allocation can assist resource sharing.

\end{enumerate}
\section{Part A: HDFS installation and verification}

Use hadoop\_\allowbreak{}installation.md and cite the exact steps in the report. Download and extract Hadoop using Step 1;\allowbreak{} configure environment variables using Step 2. Configure core-site.xml and hdfs-site.xml using Steps 3 and 4. Create dfs.namenode.name.dir and dfs.datanode.data.dir and set ownership. Set JAVA\_\allowbreak{}HOME in hadoop-env.sh following Step 5. Initialize a new laboratory NameNode with hdfs namenode -format, start the services with start-dfs.sh, and report cluster health using hdfs dfsadmin -report. Formatting is a one-time initialization of a new namespace;\allowbreak{} it must not be repeated against an existing dataset. Deliver the DataNode/\allowbreak{}capacity report screenshot and the following upload/\allowbreak{}listing test, substituting the actual user name:\allowbreak{}


\begin{quote}\small\ttfamily\raggedright
printf 'hello spark\textbackslash{}nhello hdfs\textbackslash{}nspark spark\textbackslash{}n' > sample.txt\par
hdfs dfs -mkdir -p /\allowbreak{}user/\allowbreak{}<your\_\allowbreak{}user>/\allowbreak{}lab/\allowbreak{}\par
hdfs dfs -put -f sample.txt /\allowbreak{}user/\allowbreak{}<your\_\allowbreak{}user>/\allowbreak{}lab/\allowbreak{}\par
hdfs dfs -ls /\allowbreak{}user/\allowbreak{}<your\_\allowbreak{}user>/\allowbreak{}lab/\allowbreak{}\par\end{quote}
\section{Part B: Spark Standalone and HDFS integration}

Use spark\_\allowbreak{}installation.md and identify the Install Java and Install Spark steps. Install Java, download/\allowbreak{}extract Spark, and add the environment variables to \textasciitilde{}/\allowbreak{}.bashrc. Start the master with start-master.sh and the worker with start-worker.sh spark:\allowbreak{}/\allowbreak{}/\allowbreak{}<MASTER\_\allowbreak{}IP>:\allowbreak{}7077. Inspect http:\allowbreak{}/\allowbreak{}/\allowbreak{}<MASTER\_\allowbreak{}IP>:\allowbreak{}8080 and supply a screenshot showing the registered worker. Open a Scala shell with the standalone master or local[*]. Load the HDFS test file and execute:\allowbreak{}


\begin{quote}\small\ttfamily\raggedright
val rdd =\allowbreak{} sc.textFile("hdfs:\allowbreak{}/\allowbreak{}/\allowbreak{}<NAMENODE\_\allowbreak{}IP>:\allowbreak{}9000/\allowbreak{}user/\allowbreak{}<your\_\allowbreak{}user>/\allowbreak{}lab/\allowbreak{}sample.txt")\par
val counts =\allowbreak{} rdd.flatMap(\_\allowbreak{}.split("\textbackslash{}\textbackslash{}s+\allowbreak{}"))\par
 .filter(\_\allowbreak{}.nonEmpty).map(w =\allowbreak{}> (w.toLowerCase, 1))\par
 .reduceByKey(\_\allowbreak{} +\allowbreak{} \_\allowbreak{}).sortBy \{ case (\_\allowbreak{}, c) =\allowbreak{}> -c \}\par
counts.collect().foreach(println)\par\end{quote}

For the stated sample text the correct counts are spark=\allowbreak{}3, hello=\allowbreak{}2, and hdfs=\allowbreak{}1. The assignment's example prints spark=\allowbreak{}2, which omits one occurrence on the third line. Sorting equal counts needs an explicit secondary key if deterministic tie ordering is required.

\section{Final report requirements}

Reference the guide file and section for every installation stage. Include requested command outputs and screenshots;\allowbreak{} briefly explain encountered errors and remedies;\allowbreak{} answer every conceptual question. Preserve the installation steps and troubleshooting narrative rather than presenting only a final success screenshot.

\end{document}